% Part: computability
% Chapter: recursive-functions
% Section: normal-form

\documentclass[../../../include/open-logic-section]{subfiles}

\begin{document}

\olfileid{cmp}{rec}{nft}
\olsection{স্বাভাবিক রূপের উপপাদ্য}

\begin{thm}[ক্লিনির স্বাভাবিক রূপের উপপাদ্য]
\ollabel{thm:kleene-nf}
এমন একটি আদিম পুনরাবৃত্ত সম্বন্ধ $T(e, x, s)$ এবং একটি আদিম পুনরাবৃত্ত
অপেক্ষক $U(s)$ আছে, যাদের নিচের ধর্ম রয়েছে: $f$ যেকোনো আংশিক
পুনরাবৃত্ত অপেক্ষক হলে কোনো~$e$-এর জন্য
\[
f(x) \simeq U(\umin{s}{T(e, x, s)})
\]
প্রত্যেক $x$-এর ক্ষেত্রে সত্য।
\end{thm}

\begin{explain}
স্বাভাবিক রূপের উপপাদ্যটির প্রমাণ জটিল, কিন্তু মূল ধারণাটি সরল। প্রত্যেক
আংশিক পুনরাবৃত্ত অপেক্ষকের একটি \emph{সূচক}~$e$ থাকে; স্বজ্ঞাতভাবে এটি
তার প্রোগ্রাম বা সংজ্ঞার সংকেতবাহী একটি সংখ্যা। $f(x) \fdefined$ হলে,
গণনাটিকে নিয়মমাফিক লিপিবদ্ধ করে কোনো সংখ্যা~$s$ দিয়ে সংকেতায়িত করা
যায়; আর~$s$ যে নিবেশ~$x$-এ $f$-এর গণনাটির সংকেত, তা শুধু $x$ ও
সংজ্ঞা~$e$ ব্যবহার করে আদিম পুনরাবৃত্ত উপায়ে পরীক্ষা করা যায়। ফলে
সম্বন্ধ~$T$—“সূচক~$e$-এর অপেক্ষকটির নিবেশ~$x$-এর জন্য একটি গণনা আছে,
এবং $s$ এই গণনাটির সংকেত”—আদিম পুনরাবৃত্ত। গণনার সম্পূর্ণ নথি~$s$
দেওয়া থাকলে~$s$-এর “ফল” হলো~$f(x)$-এর মান, এবং সেটিকেও~$s$ থেকে আদিম
পুনরাবৃত্ত উপায়ে পাওয়া যায়।

স্বাভাবিক রূপের উপপাদ্যটি দেখায় যে যেকোনো আংশিক পুনরাবৃত্ত অপেক্ষকের
সংজ্ঞার জন্য মাত্র একটি সীমাহীন অনুসন্ধানই প্রয়োজন। মূলত আমরা সব সংখ্যার
মধ্যে অনুসন্ধান চালিয়ে এমন একটি সংখ্যা খুঁজতে পারি, যা নিবেশ~$x$-এর
ক্ষেত্রে সূচক~$e$-এর অপেক্ষকটির কোনো গণনার সংকেত। সংখ্যাগুলি~$e$-কে
আংশিক পুনরাবৃত্ত অপেক্ষকের “নাম” হিসেবে ব্যবহার করা যায়, এবং উপপাদ্যের
সমীকরণে সংজ্ঞায়িত অপেক্ষক~$f$-কে $\cfind{e}$ লেখা যায়। লক্ষ করুন,
কোনো আংশিক পুনরাবৃত্ত অপেক্ষকের একাধিক সূচক থাকতে পারে—আসলে প্রত্যেক
আংশিক পুনরাবৃত্ত অপেক্ষকের অসীমসংখ্যক সূচক থাকে।
\end{explain}
\end{document}
